% !TeX root = ../../Book2.tex
\chapter{振荡积分}
\label{b2:ch:21}
\BookChapterNumber{b2:ch:21}
\begin{chapterabstract}
振荡积分的衰减来自相位变化造成的抵消。本章证明一维 van der Corput 引理、多维非驻相法、非退化驻相展开和混合 Hesse 矩阵 可逆时的振荡算子 \(L^2\) 估计，并解释曲率、参数退化及 Fourier 积分算子的初步几何。
\end{chapterabstract}

\begin{learninggoals}\begin{itemize}
\item\label{b2:item:21-goal-1}能区分边界项、驻点贡献与非驻相余项，并跟踪估计常数。
\item\label{b2:item:21-goal-2}能用有限阶导数下界证明一维衰减，并判断适用假设。
\item\label{b2:item:21-goal-3}能计算 Hesse 矩阵 行列式与符号，得到驻相领先项及下一阶。
\item\label{b2:item:21-goal-4}能使用 \(TT^*\) 和 Schur 估计证明局部振荡算子的 \(L^2\) 界。
\end{itemize}\end{learninggoals}

在前几章，Fourier 估计主要把振荡放进乘子或核的结构中。本章直接研究积分里的相位，说明曲率如何变成幂次衰减。所有高频参数均取 \(\lambda\ge1\)，所有相位均为实值。先处理标量积分，再进入算子估计；算子部分所需的非退化条件不同于标量驻相条件，不能互相替换。

% !TeX root = ../../Book2.tex
\section{一维振荡积分}
\label{b2:sec:2101}

正积分的大小由振幅总量决定，振荡积分还受到相位变化的制约。令
\[
 I_\lambda(a,\phi)=\int_\alpha^\beta e^{2\pi\mathrm i\lambda\phi(t)}a(t)\,\mathrm dt,
 \qquad \lambda\ge1.
\]
这里相位为实函数，振幅允许取复值。衰减估计中的常数必须说明依赖哪些导数与非退化下界；若这些量随参数恶化，单写一个与参数无关的 \(O(\lambda^{-1})\) 就失去了内容。

\subsection{分部积分}
\label{b2:subsec:210101}

假设 \(\phi'\ne0\)，并令 \(a\) 绝对连续。由
\[
 e^{2\pi\mathrm i\lambda\phi}
 =\frac1{2\pi\mathrm i\lambda\phi'}\frac{\mathrm d}{\mathrm dt}e^{2\pi\mathrm i\lambda\phi}
\]
得到边界项与余项的准确分离：
\[
 I_\lambda=\left[\frac{e^{2\pi\mathrm i\lambda\phi}a}{2\pi\mathrm i\lambda\phi'}\right]_\alpha^\beta
 -\frac1{2\pi\mathrm i\lambda}\int_\alpha^\beta e^{2\pi\mathrm i\lambda\phi}
       \left(\frac a{\phi'}\right)'\,\mathrm dt.
\]
若振幅在两端消失，第一项为零；若不消失，就必须保留它。对于线性相位与常数振幅，全部积分恰好由边界给出。这也说明增加振幅的光滑性而不处理边界，通常不能得到任意阶衰减。

\ProseParagraphBreak
在 \(|\phi'|\ge c>0\) 时，上式给出的上界为
\[
 |I_\lambda|\le\frac1{2\pi\lambda}
 \left[c^{-1}\bigl(|a(\alpha)|+|a(\beta)|+\|a'\|_1\bigr)
       +c^{-2}\int |a\phi''|\right].
\]
分母中出现的是相位导数；大的相位本身并不保证振荡，例如加上任意常数只改变整体复相位。

\subsection{非驻相}
\label{b2:subsec:210102}

\begin{proposition}[非驻相衰减]\label{b2:prop:21-nonstationary-one}
若 \(a\in C_c^\infty((\alpha,\beta))\)、\(\phi\in C^\infty\)，且在振幅支撑的一个邻域内 \(|\phi'|\ge c>0\)，则对每个整数 \(N\ge1\)，
\[
 |I_\lambda|\le C_N\lambda^{-N}.
\]
常数只依赖该邻域上 \(c^{-1}\)、相位的有限阶导数及振幅的有限阶导数。
\end{proposition}
\begin{proof}
把上一个公式反复用于振幅，定义
\(a_0=a\)、\(a_{j+1}=-(a_j/\phi')'\)。
每个 \(a_j\) 仍紧支撑于同一个集合，因此没有边界项，且
\(I_\lambda=(2\pi\mathrm i\lambda)^{-N}\int e^{2\pi\mathrm i\lambda\phi}a_N\)。
对 \(a_N\) 取 \(L^1\) 范数即得结论。商的每次求导只增加有限阶相位导数及 \(\phi'\) 的负幂，说明了常数依赖。
\end{proof}
紧支撑应理解为支撑远离正在使用公式的积分域边界。把光滑函数硬截断在一个区间上，会产生新的边界项，不能直接套用此命题。

\subsection{van der Corput 引理}
\label{b2:subsec:210103}

下面采用\term[van-der-corput-yinli]{van der Corput 引理}[van der Corput lemma]的下述版本。

\begin{theorem}[van der Corput 引理]\label{b2:thm:21-vdc}
设 \(a\) 绝对连续，\(\phi\in C^k([\alpha,\beta])\)。若 \(k=1\)，假设 \(|\phi'|\ge1\) 且 \(\phi'\) 单调；若 \(k\ge2\)，假设 \(|\phi^{(k)}|\ge1\)。则
\[
 |I_\lambda|\le C_k\lambda^{-1/k}
 \bigl(\|a\|_\infty+\|a'\|_1\bigr).
\]
\end{theorem}
\begin{proof}
先取振幅为 \(1\)。\(k=1\) 时分部积分，单调性给出
\(\int|(\phi')^{-2}\phi''|\le2\)，故任意子区间上的积分不超过 \(C\lambda^{-1}\)。
下面证明所需的次水平集估计：若 \(|h^{(r)}|\ge1\)，则
\[
 |\{t:|h(t)|\le\delta\}|\le C_r\delta^{1/r}.
\]
记左边为 \(L>0\)。按该集合的累积测度选择 \(r+1\) 个点，使任意两相邻点的距离至少 \(L/(2r)\)；可先以略低于 \(L\) 的测度作选择，再取极限。插值多项式的最高次系数为
\(\sum_{j=0}^r h(t_j)/\prod_{\ell\ne j}(t_j-t_\ell)\)，绝对值至多 \(C_r\delta L^{-r}\)。
逐次应用 Rolle 定理，某点的 \(h^{(r)}/r!\) 等于这个系数，得到所需估计。

对 \(h=\phi'\)、\(r=k-1\) 应用它。集合 \(|\phi'|\le\delta\) 上的积分至多 \(C_k\delta^{1/(k-1)}\)。其补集可分成有限个 \(\phi'\) 单调且 \(|\phi'|>\delta\) 的区间：\(\phi^{(k)}\) 无零点，Rolle 定理说明 \(\phi''\) 至多有 \(k-2\) 个零点；每个单调区间再按 \(\pm\delta\) 切分。第一阶估计在这些区间给出总上界 \(C_k(\lambda\delta)^{-1}\)。
取 \(\delta=\lambda^{-(k-1)/k}\)，两项均为 \(C_k\lambda^{-1/k}\)。

最后设 \(F(t)=\int_\alpha^t e^{2\pi\mathrm i\lambda\phi(s)}\,\mathrm ds\)。
上面的证明适用于每个子区间，所以 \(\|F\|_\infty\le C_k\lambda^{-1/k}\)。
公式 \(I_\lambda=a(\beta)F(\beta)-\int a'F\) 给出一般振幅的结论。
\end{proof}
第一阶中单调性不能删除。相位导数始终非零但反复变化时，分部积分会积累导数的总变差；只有下界而没有该变差控制，便不能保证统一常数。

\subsection{衰减估计}
\label{b2:subsec:210104}

若相位的第 \(k\) 阶导数下界为 \(c\)，对 \(\lambda c\) 应用定理，可得到 \(C_k(\lambda c)^{-1/k}\) 的上界。因此相位趋于退化时，常数会变大。以
\(\phi(t)=t^k\)、振幅在原点等于 \(1\) 为例，令 \(s=\lambda^{1/k}t\)，便出现自然尺度 \(\lambda^{-1/k}\)。对于二次相位，该尺度确实达到；更高阶时适当选取单侧振幅也可避免两侧抵消。

\ProseParagraphBreak
常数下界与有限类型条件也可局部使用。若紧支撑上每一点至少有某个 \(1\le j\le k\) 满足 \(|\phi^{(j)}|\ge c\)，有限覆盖与光滑分割把积分分成若干项。第一阶部分还需细分成有受控变差的区间；其余部分可直接使用 van der Corput 引理。这里有限覆盖的大小、分割函数导数和非退化下界都进入常数。有限类型给出一个统一衰减率，并不说明每个点都以最坏的 \(k\) 阶速度贡献。

% 扩写定位：b2:subsec:210104
有限型估计在临界点碰撞时比逐点驻相更稳定。取
\[
 I(\lambda,s)=\int e^{2\pi\mathrm i\lambda(t^3-st)}a(t)\,\mathrm dt,
 \qquad a\in C_c^\infty(\mathbb R).
\]
三阶导数恒为六，因此 van der Corput 给出
\(|I(\lambda,s)|\leq C_a\lambda^{-1/3}\)，常数与 \(s\) 无关。若 \(s>0\)，两个临界点为 \(\pm\sqrt{s/3}\)，二阶导数大小为 \(2\sqrt{3s}\)；分别使用驻相会出现 \(s^{-1/4}\lambda^{-1/2}\) 的因子，不能在 \(s\downarrow0\) 时保持统一。

\ProseParagraphBreak

换元 \(t=\lambda^{-1/3}u\) 后，相位成为 \(u^3-(s\lambda^{2/3})u\)。因此碰撞区域的自然参数是 \(s\lambda^{2/3}\)，自然宽度是 \(s\asymp\lambda^{-2/3}\)。在这个区域，两点驻相的分离距离已经与波动尺度相当，应保持三次相位整体；当 \(s\lambda^{2/3}\) 很大时，才适合分别围绕两个临界点分析。这里的统一估计无需写出 Airy 函数，但缩放已经解释了为什么二次渐近在碰撞附近失效。

\ProseParagraphBreak

非驻相区域也需要与支撑相配合。若 \(s<0\)，则 \(\phi_s'(t)=3t^2+|s|\geq|s|\)，固定的负 \(s\) 给出任意阶衰减，但分部积分常数会随 \(|s|\downarrow0\) 增长；统一三阶界仍然有效。若振幅支撑远离零点，则即使 \(s\) 接近零，一阶导数也有统一下界，此时可以恢复快速衰减。由此可见，退化不是相位公式本身的孤立属性，必须连同振幅支撑和参数范围一起判断。

% !TeX root = ../../Book2.tex
\section{驻相法}
\label{b2:sec:2102}

一阶导数消失后，分部积分的分母发生奇异，主要贡献便集中在驻点附近。非退化驻点的尺度是 \(\lambda^{-1/2}\)，其领先系数由 Hesse 矩阵 决定。以下把 Fourier 指数中的 \(2\pi\) 保留在相位内，以便与全卷的变换约定一致。

\subsection{非退化临界点}
\label{b2:subsec:210201}

\begin{definition}\label{b2:def:21-critical}
若 \(\nabla\phi(x_0)=0\)，称 \(x_0\) 为驻点；若同时
\(\det D^2\phi(x_0)\ne0\)，称其为\term[fei-tuihua-zhudian]{非退化驻点}[nondegenerate critical point]。
\end{definition}
逆函数定理用于映射 \(\nabla\phi\)，说明非退化驻点在一个小邻域内是唯一的。反过来，孤立驻点未必非退化：\(t^4\) 的原点就是例子。第二种情形的衰减阶由更高阶项决定，不能套用二次驻相公式。

\ProseParagraphBreak
在一个紧集上，若所有驻点非退化，则驻点只有有限个。否则驻点列有聚点，连续性使聚点仍为驻点，与局部孤立性矛盾。于是可以分别计算各驻点贡献，再对剩余部分使用非驻相法；不同贡献的复相位可能抵消，因此各项上界之和只是上界。

\subsection{Hesse 矩阵}
\label{b2:subsec:210202}

对实对称可逆矩阵 \(H\)，记其正特征值数减负特征值数为 \(\operatorname{sgn}H\)。在非退化驻点附近存在光滑坐标 \(u\)，使
\[
 \phi(x)=\phi(x_0)+\tfrac12\sum_{j=1}^d\varepsilon_j u_j^2,
 \qquad \varepsilon_j\in\{1,-1\}.
\]
这是局部 Morse 归约。为说明它所用的事实，先平移驻点并作固定线性变换，使 Hesse 矩阵 对角化。Taylor 积分余项写成
\(\phi(x)-\phi(0)=\tfrac12x^{\mathsf T}H(x)x\)，其中
\(H(x)=2\int_0^1(1-t)D^2\phi(tx)\,\mathrm dt\) 光滑，且 \(H(0)\) 可逆。
在足够小的邻域内，按固定顺序作对称消元，各主元仍不为零。由消元得到光滑的 \(L(x)D(x)L(x)^{\mathsf T}\) 分解；各主元的符号不变，把其绝对值的平方根吸收到三角矩阵内，便有
\(H(x)=B(x)^{\mathsf T}S B(x)\)，\(S=\operatorname{diag}(\varepsilon_j)\)。
令 \(u=B(x)x\)，其在原点的导数可逆，逆函数定理完成归约。

\ProseParagraphBreak
坐标变换的 Jacobi 因子满足
\[
 \left|\det\frac{\partial x}{\partial u}(0)\right|
      =|\det D^2\phi(x_0)|^{-1/2}.
\]
这个因子不是可忽略的规范化。一个 Hesse 矩阵 特征值变小时，对应方向的驻点邻域变宽，领先系数也随之增大。

\subsection{驻相引理}
\label{b2:subsec:210203}

下面采用\term[zhuxiang-gongshi]{驻相公式}[stationary phase formula]的下述版本。

\begin{theorem}[驻相公式]\label{b2:thm:21-stationary-phase}
设 \(a\in C_c^\infty(\mathbb R^d)\)、\(\phi\) 在支撑附近光滑，支撑内只有一个驻点 \(x_0\)，且该点非退化。记 \(H=D^2\phi(x_0)\)，则
\[
 \int e^{2\pi\mathrm i\lambda\phi(x)}a(x)\,\mathrm dx
 =e^{2\pi\mathrm i\lambda\phi(x_0)}
   \lambda^{-d/2}e^{\mathrm i\pi\operatorname{sgn}H/4}
   \frac{a(x_0)}{|\det H|^{1/2}}
   +O(\lambda^{-d/2-1}).
\]
\end{theorem}
\begin{proof}
分割振幅使一部分支撑于 Morse 坐标邻域，剩余部分由非驻相法降到任意阶。换坐标后只需处理
\(\int e^{\pi\mathrm i\lambda u^{\mathsf T}Su}b(u)\,\mathrm du\)。
对角矩阵 \(S\) 的每个分量应用复 Gaussian 积分，先加因子 \(e^{-\epsilon|u|^2}\)，再令 \(\epsilon\downarrow0\)，得到分布恒等式
\[
 \int e^{\pi\mathrm i\lambda u^{\mathsf T}Su}b(u)\,\mathrm du
 =\lambda^{-d/2}e^{\mathrm i\pi\operatorname{sgn}S/4}
   \int e^{-\pi\mathrm i\eta^{\mathsf T}S\eta/\lambda}\widehat b(\eta)\,\mathrm d\eta.
\]
Gaussian 积分在 \(\epsilon>0\) 时绝对收敛；配方所得 Fourier 变换的模有统一上界，乘上 Schwartz 函数 \(\widehat b\) 后可用支配收敛。这说明取极限与配对都是合法的。
右侧指数等于 \(1+O(|\eta|^2/\lambda)\)，而
\(\int\widehat b=b(0)\)，\(\int|\eta|^2|\widehat b|<\infty\)。
再代入 \(b(0)=a(x_0)|\det H|^{-1/2}\)，结论成立。
\end{proof}
该公式是固定振幅、固定相位下的渐近式。若振幅支撑随 \(\lambda\) 缩小或导数增长，余项常数可能增长；此时应先把自然尺度规范化。

% 扩写定位：b2:subsec:210203
边界还会改变驻相的首项。考虑一个临界点恰好位于积分端点的模型
\[
 J(\lambda)=\int_0^\infty e^{\pi\mathrm i\lambda t^2}a(t)\,\mathrm dt,
 \qquad a\in C_c^\infty([0,\infty)).
\]
此处光滑是指可延拓到零点邻域的光滑函数。若 \(a(0)\ne0\)，则
\[
 J(\lambda)=\frac12e^{\pi\mathrm i/4}a(0)\lambda^{-1/2}+O(\lambda^{-1}).
\]
取在零点附近为一的偶截断 \(\chi\)，把 \(a=a(0)\chi+t b\) 分解。第一项是全直线二次 Gaussian 积分的一半，其远离零点部分可反复分部积分；第二项令 \(v=t^2\)，成为 \(\frac12\int_0^\infty e^{\pi\mathrm i\lambda v}b(\sqrt v)\,\mathrm dv\)。该振幅的导数在零点至多有 \(v^{-1/2}\) 奇性，仍可积，分部积分得到 \(O(\lambda^{-1})\)。这就证明了首项中的一半。

\ProseParagraphBreak

若临界点离开端点而进入内部，其首项重新是完整的二次驻相系数；若移到区间外，则固定参数下没有临界点而只剩边界项。随着临界点在距离 \(\lambda^{-1/2}\) 的区域内移动，应使用保留端点的 Fresnel 积分，不能把内部公式和非驻相公式硬拼成一个统一展开。这个模型同时说明，振幅紧支于开集和振幅在闭区间端点非零是不同假设。

\subsection{渐近展开}
\label{b2:subsec:210204}

上一个证明还给出任意阶展开。令 \(D_S=\sum_j\varepsilon_j\partial_{u_j}^2\)，则
\[
 \int e^{\pi\mathrm i\lambda u^{\mathsf T}Su}b(u)\,\mathrm du
 =\lambda^{-d/2}e^{\mathrm i\pi\operatorname{sgn}S/4}
 \left[\sum_{j=0}^{N-1}\frac{(\mathrm i/(4\pi\lambda))^j}{j!}
              D_S^j b(0)+O(\lambda^{-N})\right].
\]
证明是在指数的 Taylor 余项上使用
\(|e^{\mathrm it}-\sum_{j<N}(\mathrm it)^j/j!|\le C_N|t|^N\)，再利用 Fourier 求导恒等式
\(\int(\eta^{\mathsf T}S\eta)^j\widehat b
 =(-4\pi^2)^{-j}D_S^j b(0)\)。
因此第 \(j\) 个系数由驻点处至多 \(2j\) 阶的振幅导数及坐标变换导数决定。

\ProseParagraphBreak
领先系数为零并不意味着整个积分为零。例如一维二次相位配以偶振幅 \(a(t)=t^2\chi(t)\)，首项消失，第二项通常非零，衰减变成 \(\lambda^{-3/2}\)。若振幅为奇函数、截断也对称，积分则确实为零。前者是局部消失阶带来的增益，后者是全局对称性抵消；两种机制应分清。

% 扩写定位：b2:subsec:210204
驻相公式也适用于一个临界子流形，前提是法向二次型非退化。这里可以给出一个明确的局部版本，无需提前使用 Fourier 积分算子理论。
\begin{proposition}\label{b2:prop:21-clean-stationary}
设局部坐标为 \((y,z)\in\mathbb R^{d-r}\times\mathbb R^r\)，临界集为 \(z=0\)。假定在振幅的一个固定紧支撑附近有
\(\mathrm d\phi(y,0)=0\)，法向 Hesse 矩阵 \(H(y)=\partial_z^2\phi(y,0)\) 一致可逆，且无其他临界点。若该片临界集连通，则 \(\phi(y,0)=c\) 为常数。对 \(a\in C_c^\infty\)，
\[
 \int e^{2\pi\mathrm i\lambda\phi(y,z)}a(y,z)\,\mathrm dy\,\mathrm dz
 =e^{2\pi\mathrm i\lambda c}\lambda^{-r/2}
 \int e^{\pi\mathrm i\operatorname{sgn}H(y)/4}
 \frac{a(y,0)}{|\det H(y)|^{1/2}}\,\mathrm dy
 +O(\lambda^{-r/2-1}).
\]
\end{proposition}
\begin{proof}
临界条件使沿 \(y\) 方向的一阶导数为零，所以相位在连通临界片上常值。对每个固定 \(y\)，在法向 \(z\) 上作 Taylor 展开和前面构造的 Morse 坐标变换。Hesse 矩阵 的一致可逆性与紧支撑允许在有限个参数邻域中统一选择坐标；矩阵分解、逆函数定理和振幅导数都具有统一的有限阶界。在这些片上，法向驻相公式的余项一致为 \(O(\lambda^{-r/2-1})\)。对 \(y\) 积分并加上有限分割便得结论。远离 \(z=0\) 的部分梯度具有正下界，反复分部积分把它放入任意小的幂次余项中。
\end{proof}
衰减幂次由真正振荡的法向维数 \(r\) 决定，而不是积分空间的总维数。例如 \(\phi(y,z)=|z|^2/2\) 与乘积振幅 \(a=b(y)c(z)\)，对 \(y\) 的积分没有任何 \(\lambda\) 衰减；只有 \(z\) 变量贡献 \(\lambda^{-r/2}\)。圆柱曲面的 Fourier 变换也体现这种秩信息：弯曲方向提供驻相增益，平直方向是否衰减取决于振幅和 Fourier 方向，不能从总维数套用球面的公式。

\ProseParagraphBreak

这一局部公式不是任意退化临界点的公式。若 Hesse 矩阵 在法向也退化，Morse 坐标不再存在；若临界集有交叉或尖点，也不能选成单一光滑片。比如 \(t^4\) 在零点的衰减为 \(\lambda^{-1/4}\)，已与任何一维非退化二次型的 \(\lambda^{-1/2}\) 不同。所需的是有限型分析，而不是把行列式的零值放进分母。

\ProseParagraphBreak

非二次相位的下一项并不只取决于振幅的二阶导数。这个差别可以用一维局部计算明确看到。将驻点移到零，并归一化 \(\phi''(0)=1\)，写
\[
 \phi(t)=\frac{t^2}2+\frac c6t^3+\frac d{24}t^4+O(t^5).
\]
Morse 变量为 \(y=t\sqrt{2\phi(t)/t^2}\)。在零点取使导数为正的一支，逐项展开再反解得到
\[
 t=y-\frac c6y^2+\left(\frac{5c^2}{72}-\frac d{24}\right)y^3+O(y^4).
\]
新振幅是 \(b(y)=a(t(y))t'(y)\)，而不只是 \(a(t(y))\)。计算其二阶导数得到
\[
 b''(0)=a''(0)-c\,a'(0)
 +\left(\frac{5c^2}{12}-\frac d4\right)a(0).
\]
二次模型的展开于是给出
\[
 I_\lambda=e^{\pi\mathrm i/4}\lambda^{-1/2}
 \left[a(0)+\frac{\mathrm i}{4\pi\lambda}b''(0)
 +O(\lambda^{-2})\right],
\]
这里假设相位在该局部片只有这个驻点，远离它的部分已用非驻相法处理。

\ProseParagraphBreak

即使原振幅在驻点附近为常数，三次和四次相位也会改变次项，因为坐标变换的体积因子不是常数。若只把二次公式中的 \(a''(0)\) 原样用于任意相位，就会遗漏这些项。计算首项时只需变换在驻点的一阶微分；计算更高项则需要相位与变换的更高阶 Taylor 数据。这解释了为什么完整渐近展开的导数假设随所需阶数增加。

\ProseParagraphBreak

参数族的微分也应先去掉迅速振荡的临界值。设相位和振幅光滑依赖参数 \(v\)，唯一非退化临界点 \(x_0(v)\) 位于固定片内，并有统一下界与支撑。隐函数定理使 \(x_0(v)\) 光滑。把积分乘以
\(e^{-2\pi\mathrm i\lambda\phi(x_0(v),v)}\lambda^{d/2}\)，则该归一化量具有以 \(\lambda^{-1}\) 为幂的渐近展开，系数光滑依赖 \(v\)，并在紧参数集上连同每个固定阶参数导数一致成立。

\ProseParagraphBreak

证明仍用含参数的 Morse 坐标：矩阵分解和逆映射都光滑依赖参数，新振幅的所需导数因此一致有界。二次 Fourier 计算的 Taylor 余项可以在这些半范数下统一估计。对远离临界点的部分，参数微分可能带来若干 \(\lambda\) 因子，但预先增加相应次数的非驻相分部积分即可吸收。因此“对参数一致的函数值余项”与“对参数一致的全部所需导数余项”要分别核对；只有前者不足以支持后续对参数求导。

\ProseParagraphBreak

有限个驻点的贡献需要先相加，才能讨论积分的整体大小。假定支撑中只有非退化驻点 \(x_1,\ldots,x_\ell\)，选各点附近互不相交的截断，再把其余部分留给非驻相法，便有
\[
 I_\lambda=\lambda^{-d/2}\sum_{\nu=1}^{\ell}
 e^{2\pi\mathrm i\lambda\phi(x_\nu)}
 e^{\pi\mathrm i\operatorname{sgn}H_\nu/4}
 \frac{a(x_\nu)}{|\det H_\nu|^{1/2}}
 +O(\lambda^{-d/2-1}).
\]
这里每个局部首项都可能非零，但它们的和仍可能抵消。取一维相位 \(\phi(t)=t^3/3-t\)，振幅只支在 \(\pm1\) 的两个小邻域，并在两个驻点等于一。两点的相位值为 \(\mp2/3\)，二阶导数为 \(\pm2\)，故首项合成为
\[
 \sqrt2\,\lambda^{-1/2}
 \cos\left(\frac{4\pi\lambda}{3}-\frac\pi4\right).
\]
沿余弦消失的一列正参数，积分只剩 \(O(\lambda^{-3/2})\)。因此一般驻相定理提供上界与展开；若要求统一的下界，还必须排除不同驻点之间的干涉。反过来，只有一个驻点且该点振幅非零时，首项模有正下界，余项便保证充分大参数下确实具有 \(\lambda^{-d/2}\) 的大小。

\ProseParagraphBreak

首项中的行列式与体积因子也有坐标无关的含义。设 \(x=F(y)\) 是局部光滑换元，\(F(y_0)=x_0\)，且 \(B=DF(y_0)\) 可逆。在驻点处一阶相位导数为零，二阶链式法则中的变换二阶导数项消失，因而新 Hesse 矩阵 为 \(B^{\mathsf T}HB\)，新振幅为 \(a(x_0)|\det B|\)。于是
\[
 \frac{a(x_0)|\det B|}
 {|\det(B^{\mathsf T}HB)|^{1/2}}
 =\frac{a(x_0)}{|\det H|^{1/2}}.
\]
实可逆合同也保持正负惯性指数，所以相位因子不变。这说明驻相系数与选择哪一套 Morse 坐标无关。若只改变相位而忘记变换振幅中的 Jacobi 因子，就会破坏这种不变性；前面的一维次项计算正是这一原则在更高 Taylor 阶上的体现。

% !TeX root = ../../Book2.tex
\section{多维振荡积分}
\label{b2:sec:2103}

多维情况下既可能完全没有驻点，也可能只在某些方向振荡。估计的有效维数取决于可利用的非退化方向数，而不单由积分变量的个数决定。

\subsection{相位函数}
\label{b2:subsec:210301}

若 \(\nabla\phi\ne0\)，定义
\[
 L=\frac{\nabla\phi\cdot\nabla}{2\pi\mathrm i\lambda|\nabla\phi|^2},
 \qquad Le^{2\pi\mathrm i\lambda\phi}=e^{2\pi\mathrm i\lambda\phi}.
\]
对紧支撑振幅反复转移 \(L\) 的形式伴随，可得任意阶 \(\lambda^{-N}\) 衰减。
展开伴随时出现的每一项都是相位有限阶导数、\(|\nabla\phi|^{-1}\) 的幂与振幅导数的乘积。在 \(|\nabla\phi|\ge c\) 的紧集上这些系数有界，故证明和一维完全相同，只是导数变成多重指标。

\ProseParagraphBreak
若积分域有边界，散度定理同时产生边界积分。比如平直立方体上的线性相位，即使内域完全无驻点，某些面上的切向梯度仍可能为零。这解释了为什么光滑截断与特征函数截断的 Fourier 衰减差别很大。

\subsection{振幅}
\label{b2:subsec:210302}

振幅支撑的大小、光滑性和消失阶都参与估计。给定 \(a_\lambda(x)=a(\lambda^\theta x)\)，其第 \(j\) 阶导数有 \(\lambda^{j\theta}\) 的增长；直接套用固定振幅定理会把这些增长藏进余项。先换变量 \(u=\lambda^\theta x\)，新的相位参数变为 \(\lambda^{1-2\theta}\)（二次相位），才看得清有效振荡尺度。

\ProseParagraphBreak
振幅可以包含较慢的参数，但必须统一控制其支撑和导数。一个常用情形是 \(a(x,z)\) 支撑于固定紧集，\(z\) 属于紧参数集，并且所需 \(x\) 导数一致有界。若相位也具有统一非退化下界，所有前述估计都可对 \(z\) 一致成立。参数集不紧时应另行核对这些界。

\subsection{曲率}
\label{b2:subsec:210303}

设光滑超曲面在局部写成 \(x_d=h(x')\)。曲面测度的 Fourier 变换在方向 \(\omega\) 上具有相位
\(\phi_\omega(x')=-\omega'\cdot x'-\omega_d h(x')\)。
驻点条件说明 \(\omega\) 平行于曲面的法向；驻点处 Hesse 矩阵 为
\(-\omega_dD^2h\)。在相应坐标中 Gaussian 曲率非零，等价于这个 Hesse 矩阵 可逆。因此紧支撑光滑曲面密度满足
\[
 |\widehat{a\,\mathrm d\sigma}(\xi)|\le C(1+|\xi|)^{-(d-1)/2}
\]
（假设支撑上 Gaussian 曲率处处非零）。
对方向单位球面作有限覆盖，分出驻点邻域与非驻点区域，驻相公式和非驻相法分别控制它们，便得到方向一致的常数。

\ProseParagraphBreak
若曲面是平面，沿法向相位是常数，便没有这样的衰减。若曲率只有 \(r\) 个非零方向，局部二次模型只有 \(r\) 个振荡变量，通常只能期待 \(|\xi|^{-r/2}\)。柱面是一个直接例子：圆周方向提供曲率，轴向仍可保持相干。仅说“曲面光滑”不足以推出球面型衰减。

\subsection{局部化}
\label{b2:subsec:210304}

局部化需要同时记录所选覆盖和每片的非退化下界。若 \(\chi_j\) 构成振幅支撑附近的有限光滑分割，则
\[
 I_\lambda(a,\phi)=\sum_j I_\lambda(a\chi_j,\phi).
\]
在非驻相片使用梯度下界，在驻相片使用 Hesse 矩阵 下界；余项的阶数可以统一选取，再将有限项相加。
对于一个逼近退化曲面的参数族，片数或下界可能变化。此时统一估计需要新的尺度分解，不能仅把每个固定参数的驻相公式相加。

\ProseParagraphBreak
以 \(\phi_s(t)=t^3-st\) 为例，\(s>0\) 时有两个非退化驻点，\(s=0\) 时合并为三阶驻点，\(s<0\) 时没有驻点。两个驻点的 Hesse 矩阵 在 \(s\downarrow0\) 时趋于零，分别使用二次公式的常数会发散。令 \(t=\lambda^{-1/3}u\)，可见转变发生在 \(s\lambda^{2/3}\) 为常数量级时。这个例子说明统一渐近式为何需要将参数与空间尺度一起组织。

% !TeX root = ../../Book2.tex
\section{振荡积分算子}
\label{b2:sec:2104}

标量积分的衰减不自动给出算子范数衰减，因为输入函数也会随参数改变，并可抵消一部分振荡。算子理论通过比较两个输出点，使不能同时抵消的相位差显现出来。

\subsection{算子核}
\label{b2:subsec:210401}

考虑
\[
 T_\lambda f(x)=\int_{\mathbb R^d}e^{2\pi\mathrm i\lambda\Phi(x,y)}a(x,y)f(y)\,\mathrm dy,
 \qquad a\in C_c^\infty(\mathbb R^{2d}).
\]
固定 \(x\) 后对核应用驻相法，只适用于足够光滑且导数受控的 \(f\)。\(L^2\) 单位球中的函数不满足这种控制，所以需要另一个方法。
若 \(\Phi(x,y)=\psi(x)+\theta(y)\)，则输入可取
\(f(y)=e^{-2\pi\mathrm i\lambda\theta(y)}g(y)\)，完全消去 \(y\) 振荡；输出只是固定非振荡积分乘一个模为 \(1\) 的因子，范数不衰减。

\ProseParagraphBreak
混合 Hesse 矩阵 \(\partial_x\partial_y\Phi\) 正是排除这种分离相位的量。它描述改变输出点时，输入变量中的相位梯度如何改变。对于双线性相位 \(x\cdot y\)，算子等于缩放后的 Fourier 变换，Plancherel 给出准确的 \(\lambda^{-d/2}\) 范数。

\subsection{\texorpdfstring{\(TT^*\)}{TT*} 方法}
\label{b2:subsec:210402}

伴随为
\[
 T_\lambda^*g(y)=\int e^{-2\pi\mathrm i\lambda\Phi(x,y)}
                      \overline{a(x,y)}g(x)\,\mathrm dx,
\]
故 \(T_\lambda T_\lambda^*\) 的核是
\[
 K_\lambda(x,z)=\int e^{2\pi\mathrm i\lambda(\Phi(x,y)-\Phi(z,y))}
                          a(x,y)\overline{a(z,y)}\,\mathrm dy.
\]
Hilbert 空间恒等式 \(\|T_\lambda\|_{2\to2}^2
=\|T_\lambda T_\lambda^*\|_{2\to2}\) 可由
\(\|T^*g\|^2=\langle TT^*g,g\rangle\) 与 \(\|T\|=\|T^*\|\) 直接证明。
此时核的主对角线 \(x=z\) 不振荡，但离开对角线后相位差可有统一梯度下界。估计目标因此是一个宽度约 \(\lambda^{-1}\) 的对角邻域。

\ProseParagraphBreak
若 \(\sup_x\int|K(x,z)|\,\mathrm dz\le A\) 且
\(\sup_z\int|K(x,z)|\,\mathrm dx\le B\)，则 Schur 估计给出
\(\|K\|_{2\to2}\le\sqrt{AB}\)。证明是先对
\(\int |K(x,z)||g(z)|\,\mathrm dz\) 用带权 Cauchy–Schwarz，再交换非负积分。这里绝对值估计发生在已经利用振荡得到衰减之后；若一开始对原核取绝对值，振荡增益便消失了。

\subsection{\texorpdfstring{\(L^2\)}{L2} 估计}
\label{b2:subsec:210403}

\begin{theorem}[非退化振荡算子的局部 \(L^2\) 估计]\label{b2:thm:21-oscillatory-l2}
若 \(\det\partial_x\partial_y\Phi(x_0,y_0)\ne0\)，则存在该点的邻域，使支撑在此邻域内的光滑振幅满足
\[
 \|T_\lambda f\|_2\le C\lambda^{-d/2}\|f\|_2,\qquad \lambda\ge1.
\]
若混合 Hesse 矩阵 在一个固定紧支撑附近处处可逆，同样的结论成立。
\end{theorem}
\begin{proof}
先取支撑足够小，使混合 Hesse 矩阵 一致接近固定可逆矩阵 \(H_0\)。由均值公式，
\[
 \nabla_y\Phi(x,y)-\nabla_y\Phi(z,y)
 =\left[\int_0^1\partial_x\nabla_y\Phi(z+t(x-z),y)\,\mathrm dt\right](x-z),
\]
其长度至少 \(c|x-z|\)。把相位差除以 \(|x-z|\) 后，该相位的有限阶 \(y\) 导数一致有界，梯度一致有正下界。多维非驻相估计给出
\[
 |K_\lambda(x,z)|\le C_N(1+\lambda|x-z|)^{-N}.
\]
当 \(\lambda|x-z|\le1\) 时用振幅的积分上界补上这个式子即可。取 \(N>d\)，两侧积分均不超过 \(C\lambda^{-d}\)，Schur 估计和 \(TT^*\) 恒等式得出结论。
一般紧支撑由有限个小矩形覆盖，以光滑分割写成有限个局部算子。每个局部算子满足相同幂次的估计，三角不等式给出整体结论；只改变常数。
\end{proof}
这个局部化步骤很必要：处处可逆的矩阵沿一段线段积分后未必仍可逆。证明中的固定矩阵近似只在小片内成立，不能在大支撑上直接使用。

\subsection{参数依赖性}
\label{b2:subsec:210404}

在一个参数族中，若混合 Hesse 矩阵 最小奇异值一致大于 \(c>0\)，相位与振幅的所需导数一致有界，支撑处于同一紧集，则有限覆盖可统一选择，上述常数也统一。若非退化下界趋于零，必须显示其依赖。
例如 \(\Phi_\epsilon(x,y)=\epsilon x\cdot y\)，在有效振荡范围 \(\lambda\epsilon\ge1\) 中，缩放 Fourier 变换给出 \((\lambda\epsilon)^{-d/2}\)；当 \(\lambda\epsilon<1\) 时，振幅的普通积分上界更合适。

\ProseParagraphBreak
若两侧维数不同，混合 Hesse 矩阵 不再是方阵，可逆条件也失去意义。此时应考虑其秩以及余下变量中的几何关系，不能沿用 \(\lambda^{-d/2}\) 的公式。例如把上述算子乘一个无振荡的额外变量积分，新增维度并不会增加衰减。这种秩与曲率的配合将在限制问题和多线性问题中再次出现。

% 扩写定位：b2:subsec:210404
统一常数在使用 \(TT^*\) 时有两个来源。第一是局部相位梯度差的下界，需要混合矩阵与某个固定可逆矩阵足够接近；第二是分部积分中各个振幅和倒梯度系数的导数界。仅给出行列式不为零还不够得到一族算子的统一常数，例如矩阵 \(\operatorname{diag}(\epsilon,\epsilon^{-1})\) 的行列式恒为一，但导数大小和最小奇异值都失去控制。

\ProseParagraphBreak

在有限阶估计中，实际上只需相位和振幅的有限多个导数；取 \(TT^*\) 的非驻相次数 \(N>d\) 后，所需阶数由 \(N,d\) 决定。若要全部阶的渐近展开，则应有全部相应阶导数。两种命题虽然都常写为“光滑”，对参数族的要求并不相同。记录常数时应说明固定了哪些半范数，而不能把全部依赖藏在一个无法随参数核对的 \(C\) 中。

\ProseParagraphBreak

该 \(\lambda^{-d/2}\) 幂次在一般非退化算子中不能改善。取 \(\Phi(x,y)=x\cdot y\)、振幅 \(a(x,y)=\chi(x)\eta(y)\)，其中 \(\chi,\eta\) 在零点附近不为零。固定 \(g\in C_c^\infty\) 满足 \(\eta g\not\equiv0\)，令输入为 \(g\)。则
\[
 T_\lambda g(x)=\chi(x)\widehat{\eta g}(-\lambda x),\qquad
 \|T_\lambda g\|_2^2=
 \lambda^{-d}\int|\chi(v/\lambda)|^2|\widehat{\eta g}(-v)|^2\,\mathrm dv.
\]
支配收敛给出 \(\lambda^d\|T_\lambda g\|_2^2\to|\chi(0)|^2\|\eta g\|_2^2>0\)。因而算子范数至少是 \(c\lambda^{-d/2}\)。这个例子直接测试了上界的幂次，所用输入不必依赖于 \(\lambda\)。

% !TeX root = ../../Book2.tex
\OptionalSection{Besov 与 Triebel–Lizorkin 空间}{b2:sec:1905}

同样一组频率块，可以记录两种不同的信息。先量每一块在整个空间中的大小，再比较不同尺度，得到一个数列范数；先在每个位置比较各尺度，再对位置积分，得到另一种范数。\subsecref{b2:subsec:190404}的平方空间已经采用了后一种次序。本节允许尺度指数变化，并通过差分、乘积和平坦边界的迹，说明这个指数为什么不能省略。一般空间可以选读；\chapref{b2:ch:22}的普通乘子正文只需要上一节已经建立的平方情形。

\subsection{二进定义}
\label{b2:subsec:190501}

先把分割固定下来，使后面的每个范数都有确定含义。低频块的作用也必须保留：所有齐次环带都消去多项式，而非齐次分解还能记录零频附近的信息。

\Needspace{6\baselineskip}
\begin{definition}\label{b2:def:19-besov-triebel-lizorkin}
设整数 \(d\geq1\)，取实径向函数 \(\chi\in C_c^\infty(\mathbb R^d)\)，满足 \(0\leq\chi\leq1\)，在 \(|\xi|\leq1\) 为一、\(|\xi|\geq2\) 为零，并沿半径单调不增。令 \(\varphi=\chi-\chi(2\,\cdot)\)。对复值温和分布 \(f\in\mathcal S'(\mathbb R^d)\)，定义
\[
 P_0f=\mathcal F^{-1}(\chi\cdot\widehat f),
 \quad P_jf=\mathcal F^{-1}\bigl(\varphi(2^{-j}\,\cdot)\cdot\widehat f\bigr)
 \quad(j\geq1).
\]
设 \(s\in\mathbb R\)、\(1\leq p,q\leq\infty\)。满足
\[
 \|f\|_{B^s_{p,q}}
 \coloneqq\left\|\bigl\{2^{js}\|P_jf\|_p\bigr\}_{j\geq0}\right\|_{\ell^q}<\infty
\]
的温和分布组成\term[besov-kongjian]{Besov 空间}[Besov space]，记为 \(B^s_{p,q}(\mathbb R^d)\)。当 \(1\leq p<\infty\)、\(1\leq q\leq\infty\) 时，满足
\[
 \|f\|_{F^s_{p,q}}
 \coloneqq\left\|\left\|\bigl\{2^{js}P_jf(x)\bigr\}_{j\geq0}\right\|_{\ell^q}\right\|_{L^p_x}<\infty
\]
的温和分布记为 \(F^s_{p,q}(\mathbb R^d)\)，称为一般的 Triebel--Lizorkin 空间。序列指数 \(q=\infty\) 时，\(\ell^q\) 范数取上确界；Besov 定义中的 \(p=\infty\) 取空间本性上确界。
\end{definition}
\symidx{\(B^s_{p,q}\)：非齐次 Besov 空间}
\symidx{\(F^s_{p,q}\)：非齐次 Triebel--Lizorkin 空间}
这里的块正是\cref{b2:def:19-inhomogeneous-square}中的块。每个块都是光滑函数，但未必有有限的全空间范数；定义的有限性不是自动得到的。由于
\[
 P_0f+\sum_{j=1}^NP_jf
 =\mathcal F^{-1}\bigl(\chi(2^{-N}\,\cdot)\cdot\widehat f\bigr)
 \longrightarrow f\quad\text{于 }\mathcal S',
\]
这些范数为零时所有块为零，继而 \(f=0\)。三角不等式分别来自 \(L^p\) 与 \(\ell^q\) 的三角不等式，所以在本节的指数范围内它们确实是范数。

\ProseParagraphBreak

三个参数分别回答三个问题：\(s\) 规定高频块要乘多大的尺度权；\(p\) 规定怎样量空间大小；\(q\) 规定允许多少个尺度同时显著。只写光滑阶数而不写后两个指数，通常不能判断函数是否属于该空间。

\subsection{Sobolev 空间的重新解释}
\label{b2:subsec:190502}

\Cref{b2:prop:19-inhomogeneous-identification}已证明
\[
 F^0_{p,2}(\mathbb R^d)=L^p(\mathbb R^d)
 \quad(1<p<\infty),
 \quad F^s_{2,2}(\mathbb R^d)=H^s(\mathbb R^d)
 \quad(s\in\mathbb R),
\]
等号指空间相同、范数等价。前一识别包括反向的函数代表问题：若温和分布的非齐次平方函数属于 \(L^p\)，合成估计使其块和在 \(L^p\) 中收敛，分布重构再将极限识别为原分布。不能只对已知属于 \(L^p\) 的函数证明不等式，就宣布对所有温和分布完成了空间识别。

\ProseParagraphBreak

在 \(p=q=2\) 时，Tonelli 允许交换位置积分与尺度求和，直接给出 \(B^s_{2,2}=F^s_{2,2}=H^s\)。一般 \(p=q<\infty\) 也有定义范数完全相等的恒等式
\[
 \sum_{j\geq0}2^{jsp}\|P_jf\|_p^p
 =\int_{\mathbb R^d}\sum_{j\geq0}2^{jsp}|P_jf(x)|^p\,\mathrm dx.
\]
这一次是两个积分次序的交换；\(H^s\) 与频块范数的比较还涉及频率权重，所以通常只有等价而不是数值相等。

\ProseParagraphBreak

另一个常用描述从符号 \((1+4\pi^2|\xi|^2)^{s/2}\) 定义 \(J^s\)，再要求 \(J^sf\in L^p\)。该符号光滑、各阶导数至多作多项式增长，因而可作用于温和分布。对于 \(1<p<\infty\)，把这个条件与 \(F^s_{p,2}\) 识别，需要统一控制环带上的归一化符号及其逆；这一乘子论证见\cref{b2:prop:22-bessel-square}。本章的平方可积情形已由 Plancherel 独立完成。齐次势 \(I_\alpha\) 在零频奇异，而 \(J^{-\alpha}\) 的符号在那里有限；这种低频差别正是全空间同指数估计不同的原因，不能把二者混作同一个算子。

\subsection{Besov 空间}
\label{b2:subsec:190503}

先说明所选分割如何影响 Besov 范数。若 \((\widetilde P_k)\) 来自另一套满足定义条件的截断，则高频新块只与固定个数的旧块相交，低频处也只有固定个数的例外。分布重构因此给出有限和
\[
 \widetilde P_kf=\sum_{j\sim k}\widetilde P_kP_jf.
\]
这里 \(j\sim k\) 表示 \(|j-k|\leq L\)，可增大固定整数 \(L\) 以覆盖低频例外。新块卷积核的 \(L^1\) 范数一致有界，且在这些指标上 \(2^{ks}\) 与 \(2^{js}\) 可比，故
\[
 2^{ks}\|\widetilde P_kf\|_p
 \leq C_s\sum_{j\sim k}2^{js}\|P_jf\|_p.
\]
对序列取 \(\ell^q\) 范数，有限次平移不改变该范数，便得到一个方向的界；交换新旧分割即得反向界。这个证明覆盖 \(1\leq p,q\leq\infty\)，常数依赖两套固定截断。

\ProseParagraphBreak

为了具体看见第二指数，取 \(s>0\)，在彼此相隔固定个数二进尺度的频率小球内选 \(f_n\in\mathcal S\)，使 \(\|f_n\|_2=1\)。记相应尺度为 \(j_n=Ln\)，其中固定 \(L\) 足够大。令
\[
 f=\sum_{n\geq1}a_n2^{-j_ns}f_n,
\]
并先设 \((a_n)\) 有界。不同球互不相交，正 \(s\) 使级数在 \(L^2\) 中收敛。每个球只与固定个数的频块相交，不同球对应的指标组互不交叠；在每组中，分割和为一，至少有一个块记录了固定比例的平方大小。Plancherel 和有限个数项的范数等价因此给出
\[
 \|f\|_{B^s_{2,q}}\asymp\|(a_n)\|_{\ell^q}.
\]
于是 \(a_n=1/n\) 给出属于 \(B^s_{2,2}\) 而不属于 \(B^s_{2,1}\) 的函数；\(a_n=1\) 给出属于 \(B^s_{2,\infty}\) 而不属于 \(H^s\) 的函数。频率位置相同的尺度权并不能替代尺度之间的可和性。

\ProseParagraphBreak

当光滑阶数小于一时，还可以不用 Fourier 变换来辨认一种 Besov 范数。频率块在步长远小于其波长时变化较小，在步长较大时只能用两倍振幅控制；下面的两段几何级数正来自这一区别。

\Needspace{6\baselineskip}
\begin{proposition}\label{b2:prop:19-besov-difference}
设整数 \(d\geq1\)、\(0<s<1\)、\(1\leq p<\infty\)，采用\cref{b2:def:19-besov-triebel-lizorkin}的非齐次块。对每个 \(f\in L^p(\mathbb R^d;\mathbb C)\)，有
\[
 \|f\|_{B^s_{p,p}}\asymp
 \|f\|_p+
 \left(\int_{\mathbb R^d}
 \frac{\|f(\,\cdot+h)-f\|_p^p}{|h|^{d+sp}}\,\mathrm dh\right)^{1/p}.
\]
允许两边同时为无穷；比较常数只依赖 \(d,s,p\) 与固定截断。因此该 Besov 空间与按右侧差分范数定义的分数 Sobolev 空间相同。
\end{proposition}
\begin{proof}
先设左边有限，记 \(u_j=P_jf\)、\(a_j=2^{js}\|u_j\|_p\)。由于 \(s>0\)，Hölder 的序列形式（\(p=1\) 时直接求和）给出
\[
 \sum_{j\geq0}\|u_j\|_p
 =\sum_{j\geq0}2^{-js}a_j
 \leq C_{s,p}\|(a_j)\|_{\ell^p}.
\]
所以块和在 \(L^p\) 中绝对收敛；分布重构说明其和就是 \(f\)，并且已经得到所需的 \(\|f\|_p\) 界。

沿线段积分及 Bernstein 估计给出 \(\|u_j(\cdot+h)-u_j\|_p\leq C2^j|h|\cdot\|u_j\|_p\)；平移不变性另给出两倍函数范数的界。因此
\[
 \|u_j(\,\cdot+h)-u_j\|_p
 \leq C\min\{1,2^j|h|\}\cdot\|u_j\|_p.
\]
在 \(2^{-k-1}<|h|\leq2^{-k}\)、\(k\geq0\) 上，对绝对重构求差并相加，得到
\[
 2^{ks}\|f(\,\cdot+h)-f\|_p
 \leq C\sum_{j\leq k}2^{-(k-j)(1-s)}a_j
       +C\sum_{j>k}2^{-(j-k)s}a_j.
\]
右边是两段可和几何序列与 \((a_j)\) 的卷积，离散 Young 不等式控制其 \(\ell^p_k\) 范数。每个步长环带上的 \(\int|h|^{-d}\,\mathrm dh\) 相同，且 \(|h|^{-sp}\) 与 \(2^{ksp}\) 可比，所以这个序列界正给出 \(|h|\leq1\) 的差分积分。\(|h|>1\) 的部分由 \(2\|f\|_p\) 以及 \(\int_{|h|>1}|h|^{-d-sp}\,\mathrm dh<\infty\) 控制。

反向设右边有限。对 \(j\geq1\)，令 \(k_j(h)=2^{jd}\check\varphi(2^jh)\)。其积分为零，因此在 \(L^p\) 中有
\[
 P_jf=\int k_j(h)\cdot\bigl(f(\,\cdot-h)-f\bigr)\,\mathrm dh.
\]
Minkowski 后对测度 \(|k_j(h)|\,\mathrm dh\) 使用 Hölder，得到
\[
 \|P_jf\|_p^p
 \leq\|k_j\|_1^{p-1}
       \int|k_j(h)|\cdot\|f(\,\cdot-h)-f\|_p^p\,\mathrm dh.
\]
\(p=1\) 时同式直接成立。Schwartz 衰减又给出
\[
 \sum_{j\geq1}2^{jsp}|k_j(h)|\leq C|h|^{-d-sp}\quad(h\ne0).
\]
具体地，写 \(r=|h|\)，每项至多为 \(Cr^{-d-sp}(2^jr)^{d+sp}(1+2^jr)^{-N}\)。在 \(2^jr\leq1\) 与 \(>1\) 两侧分别求几何级数，取 \(N>d+sp\) 即可。乘以差分范数后用 Tonelli，便控制全部高频项；低频块用 \(\|P_0f\|_p\leq\|\check\chi\|_1\cdot\|f\|_p\)。两向估计由 \(L^p\) 差分与各频带的范数联系起来。
\end{proof}
\(s<1\) 保证小步长对低频块的几何衰减可和，\(s>0\) 保证高频重构与大步长积分收敛。若 \(s\geq1\)，一般必须换用更高阶差分；对非零光滑梯度，\(s=1\) 的一阶差分积分已经出现 \(\mathrm dr/r\) 的发散。

\ProseParagraphBreak

这个判别可以在一个跳跃函数上精算。取单位立方体 \(Q=(0,1)^d\)。当每个 \(|h_i|<1\) 时，交集体积为 \(\prod_i(1-|h_i|)\)，因而对任意有限 \(p\geq1\)，
\[
 \|\mathbf1_Q(\,\cdot+h)-\mathbf1_Q\|_p^p
 =|Q\mathbin{\triangle}(Q-h)|
 =2\left(1-\prod_{i=1}^d(1-|h_i|)\right).
\]
若 \(|h|<1/2\)，括号介于 \(\max_i|h_i|\) 与 \(\sum_i|h_i|\) 之间，故与 \(|h|\) 可比。小步长积分于是与 \(\int_0^{1/2}r^{-sp}\,\mathrm dr\) 同时收敛。得到精确条件
\[
 \mathbf1_Q\in B^s_{p,p}\quad\Longleftrightarrow\quad sp<1
 \quad(0<s<1,\ 1\leq p<\infty).
\]
区域内外的函数都很简单，限制光滑阶数的是跨越边界时保留的跳跃。

\ProseParagraphBreak

绝对尺度求和还控制函数乘积。由 Bernstein，
\[
 \sum_{j\geq0}\|P_jf\|_\infty
 \leq C\sum_{j\geq0}2^{jd/p}\|P_jf\|_p
 =C\|f\|_{B^{d/p}_{p,1}}.
\]
所以 \(B^{d/p}_{p,1}\) 有一个有界连续代表：每块连续，上确界范数绝对可和，故块和一致收敛；分布重构识别这个代表。这也说明为什么临界处的序列指数一有特殊作用。

\Needspace{6\baselineskip}
\begin{proposition}\label{b2:prop:19-besov-algebra}
设整数 \(d\geq1\)、\(1\leq p<\infty\)、\(s\geq d/p\)，采用\cref{b2:def:19-besov-triebel-lizorkin}的非齐次 Besov 空间。则它的元素具有有界连续代表，其逐点乘积满足
\[
 \|fg\|_{B^s_{p,1}}
 \leq C_{d,s,p,\chi}\|f\|_{B^s_{p,1}}\cdot\|g\|_{B^s_{p,1}}.
\]
特别地，\(B^s_{p,1}(\mathbb R^d)\) 在乘法下封闭。
\end{proposition}
\begin{proof}
令 \(f_k=P_kf\)、\(g_l=P_lg\)。因为 \(s\geq d/p>0\)，两组块都在 \(L^p\) 和 \(L^\infty\) 中绝对求和，前述嵌入给出有界连续代表。双级数 \(\sum_{k,l}f_kg_l\) 在 \(L^p\) 中绝对收敛，并等于 \(fg\)。

每项乘积的 Fourier 支集包含于两频率球之和，故在半径 \(C2^{\max\{k,l\}}\) 的球内。这一支集规则可以对紧支集频率分布作卷积来理解；也可先用相同频率截断将 \(L^p\) 光滑逼近保持在略大的固定频带内，再把 Schwartz 函数的卷积公式传到分布极限。因而存在固定整数 \(L\)，使 \(P_j(f_kg_l)=0\) 当 \(j>\max\{k,l\}+L\)。

若 \(k\geq l\)，由块核的统一 \(L^1\) 界及 \(s>0\)，
\[
 \begin{aligned}
 \sum_{j\geq0}2^{js}\|P_j(f_kg_l)\|_p
 &\leq C\sum_{0\leq j\leq k+L}2^{js}\|f_k\|_p\cdot\|g_l\|_\infty\\
 &\leq C_s2^{ks}\|f_k\|_p\cdot\|g_l\|_\infty.
 \end{aligned}
\]
\(l>k\) 时交换两因子。先对有限双和估计，再用 \(L^p\) 收敛及 Fatou，得到
\[
 \begin{aligned}
 \|fg\|_{B^s_{p,1}}
 \leq C_s&\left(\sum_k2^{ks}\|f_k\|_p\right)
          \left(\sum_l\|g_l\|_\infty\right)\\
 +C_s&\left(\sum_k\|f_k\|_\infty\right)
          \left(\sum_l2^{ls}\|g_l\|_p\right).
 \end{aligned}
\]
两组上确界范数和均由对应 Besov 范数控制，即得结论。证明只用了 \(s\geq d/p\) 与 \(s>0\)，所以包括 \(s=d/p\) 的临界情形。
\end{proof}
若两个相乘的块频率相差很大，输出仍以较高频为主；若二者频率接近，却可能相互抵消而产生低频。上面求和的是所有 \(j\leq\max\{k,l\}+L\)，已经把这些低频输出包括在内，不能只保留高低频相互作用。这项乘积结论依赖已得到的函数代表，不能推及任意两个温和分布。

\subsection{Triebel--Lizorkin 空间}
\label{b2:subsec:190504}

这类范数先在位置 \(x\) 处比较不同频率，再作空间积分。若许多块集中在同一小区域，点态序列范数会同时记录它们；若这些块分布在彼此远离的区域，同一位置只见到少数项。先把每块的全空间范数算成一个数，会丢掉这种位置关系。两种次序可以比较，却通常不能交换。

\Needspace{6\baselineskip}
\begin{proposition}\label{b2:prop:19-besov-triebel-embeddings}
设整数 \(d\geq1\)、\(s\in\mathbb R\)、\(1\leq p<\infty\)、\(1\leq q\leq\infty\)，所有空间使用\cref{b2:def:19-besov-triebel-lizorkin}中的同一分割。则有连续嵌入
\[
 B^s_{p,\min\{p,q\}}\subset F^s_{p,q}
 \subset B^s_{p,\max\{p,q\}}.
\]
此外，对 \(1\leq p\leq r\leq\infty\)、\(1\leq q\leq\infty\)、\(s\in\mathbb R\)，有
\[
 B^s_{p,q}\subset B^{s-d(1/p-1/r)}_{r,q},
\]
其中 \(1/\infty=0\)。
\end{proposition}
\begin{proof}
先对有限个块令 \(a_j(x)=2^{js}|P_jf(x)|\)。若 \(q\leq p\)，Minkowski 在 \(L^{p/q}\) 中给出
\[
 \left\|\sum_j a_j^q\right\|_{p/q}
 \leq\sum_j\|a_j^q\|_{p/q},
\]
即 \(\|f\|_{F^s_{p,q}}\leq\|f\|_{B^s_{p,q}}\)。逐点的 \(\ell^p\leq\ell^q\) 又给出 \(\|f\|_{B^s_{p,p}}\leq\|f\|_{F^s_{p,q}}\)。若 \(p\leq q<\infty\)，Minkowski 在序列空间 \(\ell^{q/p}\) 中给出
\[
 \left(\sum_j\left[\int a_j^p\right]^{q/p}\right)^{p/q}
 \leq\int\left(\sum_j a_j^q\right)^{p/q},
\]
从而 \(\|f\|_{B^s_{p,q}}\leq\|f\|_{F^s_{p,q}}\)；逐点的 \(\ell^q\leq\ell^p\) 给出另一边。\(q=\infty\) 时直接使用
\[
 \sup_j\|a_j\|_p\leq\|\sup_j a_j\|_p
 \leq\left(\sum_j\|a_j\|_p^p\right)^{1/p}.
\]
有限块数趋于无穷，单调收敛证明第一组嵌入。

第二组嵌入对每个块使用 Bernstein：
\[
 2^{j[s-d(1/p-1/r)]}\|P_jf\|_r
 \leq C2^{js}\|P_jf\|_p.
\]
低频块尺度为一，满足同样估计。取 \(\ell^q\) 范数即得结论，包括两个无穷端点。
\end{proof}
这些都是范数不等式；只有另有严格性例子时，才可把包含号改成真包含。\(p=q\) 时第一组嵌入就是前面由 Tonelli 得到的相等，因而不可能一概宣称严格。

\ProseParagraphBreak

分割独立性也要区分指数。对 \(q=2\)、\(1<p<\infty\)，写
\(\widetilde P_kf=\sum_{|j-k|\leq L}\widetilde P_kP_jf\)，并将有限个相邻偏移分别处理。每个固定偏移对应一族 \(\widetilde P_k\)：其缩放核有统一的标准核常数，平方范数一致有界，所以\cref{b2:prop:18-hilbert-valued-cz}控制 \(L^p(\ell^2)\) 范数。权重 \(2^{ks}\) 与 \(2^{js}\) 在相邻指标上可比，得到新旧 \(F^s_{p,2}\) 范数等价。\(p=1\) 或一般 \(q\) 不在这一 Hilbert 值强型论证的范围内。这里的定义仍使用固定分割，前面的差分、嵌入和例子也均已在声明的分割下证明；一般 Triebel--Lizorkin 分割独立性的进一步理论可参看 Hans Triebel 的《Theory of Function Spaces》\cite{Triebel1983Theory}，不能用标量 Young 不等式代替其中的向量估计。

\subsection{Besov 空间的平面迹与延拓}
\label{b2:subsec:190505}

最后考虑在 \(x_d=0\) 上取边界值。一个频率为 \(2^j\) 的块，在法向上可变化的长度约为 \(2^{-j}\)。把这一层的 \(d\) 维积分变成边界的 \((d-1)\) 维积分，函数范数会出现 \(2^{j/p}\)，因此 \(s\) 阶的自然边界阶数为 \(s-1/p\)。下面把这项尺度计算落实为一个有界迹与右逆，范围限定为平坦边界、序列指数一。

\Needspace{6\baselineskip}
\begin{proposition}[平坦边界的 Besov 迹与延拓]
\label{b2:prop:19-flat-besov-trace}
设整数 \(d\geq2\)、\(1\leq p<\infty\)、\(s>1/p\)。在 \(\mathbb R^d\) 与 \(\mathbb R^{d-1}\) 上各固定\cref{b2:def:19-besov-triebel-lizorkin}的一套非齐次分割。Schwartz 函数上的限制 \(u\mapsto u(\,\cdot,0)\) 有唯一有界线性延拓
\[
 \operatorname{Tr}:B^s_{p,1}(\mathbb R^d)
 \longrightarrow B^{s-1/p}_{p,1}(\mathbb R^{d-1}),
\]
并存在有界线性算子 \(E\) 从目标空间映回 \(B^s_{p,1}(\mathbb R^d)\)，满足 \(\operatorname{Tr}Eg=g\)。各常数只依赖固定参数和截断。
\end{proposition}
\begin{proof}
令 \(u_j=P_ju\)。其法向频率不超过 \(C2^j\)。取一维 Schwartz 函数 \(k\)，使 \(\widehat k=1\) 于 \([-C,C]\)，则由法向频率截断，
\[
 u_j(x',0)=\int_{\mathbb R}2^jk(-2^jt)\cdot u_j(x',t)\,\mathrm dt.
\]
对每个 \(x'\) 使用 Hölder，\(p'=p/(p-1)\) 给出核的 \(L^{p'}\) 范数 \(2^{j/p}\|k\|_{p'}\)；\(p=1\) 时使用 \(p'=\infty\)。再积分切向变量，得到
\[
 \|u_j(\,\cdot,0)\|_{L^p(\mathbb R^{d-1})}
 \leq C2^{j/p}\|u_j\|_{L^p(\mathbb R^d)}.
\]
卷积表示给出所用的法向连续代表；频率块本身光滑，因此这就是其普通限制。边界频率仍在半径 \(C2^j\) 的球内。记边界块为 \(P'_k\)，其核的 \(L^1\) 范数一致有界，而且 \(P'_k[u_j(\cdot,0)]=0\) 当 \(k>j+L\)。于是
\[
 \begin{aligned}
 \|u_j(\,\cdot,0)\|_{B^{s-1/p}_{p,1}}
 &\leq C\sum_{0\leq k\leq j+L}2^{k(s-1/p)}
              \cdot2^{j/p}\|u_j\|_p\\
 &\leq C_s2^{js}\|u_j\|_p.
 \end{aligned}
\]
最后的几何级数使用了 \(s-1/p>0\)。对 \(j\) 求和，边界块的限制之和在 \(L^p\) 中绝对收敛，记为 \(\operatorname{Tr}u\)；逐个边界块取极限并用 Fatou，得到目标 Besov 界。

补足延拓的唯一性。因为序列指数为一，有限频率和在 \(B^s_{p,1}\) 中稠密：对剩余块使用有限重叠及统一卷积界，其范数受原尺度序列的尾和控制。若一个有限频率和 \(v\in L^p\) 的频率支集在固定球内，选光滑紧频率乘子 \(\eta\)，在该球邻域内为一。取 \(h_n\in C_c^\infty\) 在 \(L^p\) 中逼近 \(v\)，则 \(\eta(D)h_n\in\mathcal S\)，且
\[
 \|\eta(D)h_n-v\|_p
 \leq\|\check\eta\|_1\cdot\|h_n-v\|_p\longrightarrow0.
\]
误差有共同有界频带，故只含有限个非齐次块，\(L^p\) 收敛也给出 \(B^s_{p,1}\) 收敛。由此证明 Schwartz 类稠密。对 Schwartz 函数，非齐次截断在 Schwartz 拓扑中趋于原函数，边界限制与所定义的块和相同，所以所得连续延拓唯一。

为构造右逆，将 \(g\) 分解为边界块 \(g_j=P'_jg\)。选一维 Schwartz 函数 \(\rho\)，使 \(\rho(0)=1\) 且 \(\widehat\rho\) 紧支撑，令
\[
 v_j(x',t)=g_j(x')\cdot\rho(2^jt),
 \quad Eg=\sum_{j\geq0}v_j.
\]
每项 Fourier 支集在半径 \(C2^j\) 的球内，且
\[
 \|v_j\|_{L^p(\mathbb R^d)}
 =2^{-j/p}\|g_j\|_{L^p(\mathbb R^{d-1})}\cdot\|\rho\|_{L^p(\mathbb R)}.
\]
对体内输出块再求和，由 \(s>0\) 得
\[
 \|v_j\|_{B^s_{p,1}(\mathbb R^d)}
 \leq C_s2^{js}\|v_j\|_p
 \leq C2^{j(s-1/p)}\|g_j\|_p.
\]
右侧可和，所以 \(\sum v_j\) 在 \(L^p\) 中收敛，并由块估计与 Fatou 得到 \(Eg\in B^s_{p,1}\) 及相应范数界；同一尾和估计也证明级数在 Besov 范数中收敛。每个 \(v_j\) 只有有限多个体内频率块，其块和就是本身，故由定义 \(\operatorname{Tr}v_j=v_j(\cdot,0)=g_j\)。连续性将有限和等式传给极限，得到 \(\operatorname{Tr}Eg=\sum_jg_j=g\)。
\end{proof}
临界 \(s=1/p\) 时，边界输出尺度的求和不再有几何衰减，不能由本证明得到同样的 Besov 目标；但是只求边界 \(L^p\) 范数时，绝对求和还能保留一个临界版本，见\cref{b2:ex:19-new-3}。曲边界则另需坐标变换的有界性，不能把平面上的公式直接搬到任意区域。迹的尺度损失、序列指数与边界几何，各自承担不同的条件。


\begin{chaptersummary}
非驻相区域可反复分部积分，非退化驻点提供每个方向半个幂次的衰减。van der Corput 引理把高阶退化与较慢衰减联系起来。对输入允许任意 \(L^2\) 振荡的算子，混合 Hesse 矩阵 与 \(TT^*\) 才能揭示统一增益。局部化、振幅导数、边界处理和参数下界共同决定常数；它们都是估计的一部分。
\end{chaptersummary}

\BookExercises

% 自拟：围绕本章已经建立的工具，改变条件或组合方法。
\begin{exercise}\label{b2:ex:21-new-1}
计算 \(\int_0^1e^{2\pi\mathrm i\lambda t}\,\mathrm dt\)。说明为什么这个例子不能仅凭相位无驻点就宣称任意阶衰减。
\end{exercise}
\begin{solution}
积分为 \((e^{2\pi\mathrm i\lambda}-1)/(2\pi\mathrm i\lambda)\)。沿半整数参数，其绝对值为 \(1/(\pi\lambda)\)。振幅在边界不消失，分部积分留下边界项，因此非驻相命题的紧支撑条件不满足。
\end{solution}

% 自拟：围绕本章已经建立的工具，改变条件或组合方法。
\begin{exercise}\label{b2:ex:21-new-2}
设 \(\phi''\ge c>0\)、\(a\) 绝对连续，证明 \(I_\lambda=O((\lambda c)^{-1/2})\)，并说明区间长度为什么无需进入主常数。
\end{exercise}
\begin{solution}
对相位 \(\phi/c\) 使用二阶 van der Corput 引理，参数为 \(\lambda c\)，得上界 \(C(\lambda c)^{-1/2}(\|a\|_\infty+\|a'\|_1)\)。次水平集 \(|\phi'|\le\delta\) 的长度由导数下界控制为 \(2\delta/c\)，其余区间使用单调第一阶估计，所以没有额外的区间长度项。
\end{solution}

% 自拟：围绕本章已经建立的工具，改变条件或组合方法。
\begin{exercise}\label{b2:ex:21-new-3}
对 \(a\in C_c^\infty(\mathbb R)\)，求 \(\int e^{\pi\mathrm i\lambda t^2}a(t)\,\mathrm dt\) 的前两项。
\end{exercise}
\begin{solution}
二次模型公式给出 \(\lambda^{-1/2}e^{\mathrm i\pi/4}[a(0)+\mathrm i a''(0)/(4\pi\lambda)]+O(\lambda^{-5/2})\)。若 \(a(0)=0\) 而 \(a''(0)\ne0\)，领先阶变为 \(\lambda^{-3/2}\)。
\end{solution}

% 自拟：围绕本章已经建立的工具，改变条件或组合方法。
\begin{exercise}\label{b2:ex:21-new-4}
设 \(H\) 是实对称可逆矩阵，计算相位 \(\phi(x)=x^{\mathsf T}Hx/2\) 的驻相主项，分别讨论正定与不定的情形。
\end{exercise}
\begin{solution}
主项为 \(\lambda^{-d/2}e^{\mathrm i\pi\operatorname{sgn}H/4}a(0)|\det H|^{-1/2}\)。正定时符号为 \(d\)；若有 \(r\) 个正、\(d-r\) 个负特征值，符号为 \(2r-d\)。绝对值系数相同，复相位不同。
\end{solution}

% 自拟：围绕本章已经建立的工具，改变条件或组合方法。
\begin{exercise}\label{b2:ex:21-new-5}
设 \(a\in C_c^\infty(\mathbb R^d)\)。证明相位 \(\xi\cdot x\) 的积分在 \(|\xi|\to\infty\) 时任意阶衰减，并把它与 Schwartz 函数的 Fourier 变换联系起来。
\end{exercise}
\begin{solution}
选择坐标 \(j\) 使 \(|\xi_j|\ge|\xi|/\sqrt d\)，沿该方向分部积分 \(N\) 次，得 \(|\int e^{2\pi\mathrm i\xi\cdot x}a|\le(2\pi|\xi_j|)^{-N}\|\partial_j^Na\|_1\)。对所有 \(j\) 取最大值即可。乘以任意坐标幂后仍适用同样估计，因此 Fourier 变换属于 Schwartz 空间。
\end{solution}

% 自拟：围绕本章已经建立的工具，改变条件或组合方法。
\begin{exercise}\label{b2:ex:21-new-6}
对 \(\Phi_\epsilon(x,y)=\epsilon x\cdot y\)，解释 \(\epsilon\downarrow0\) 时振荡算子估计为何不可能有统一的 \(\lambda^{-d/2}\) 常数。
\end{exercise}
\begin{solution}
取分离振幅 \(a(x,y)=\chi(x)\chi(y)\) 及非零固定非负输入 \(f\)。当 \(\lambda\epsilon\to0\)，支配收敛使输出趋于 \(\chi(x)\int\chi f\)，其范数为正。若取 \(\epsilon=\lambda^{-2}\)，统一的 \(\lambda^{-d/2}\) 界将使该范数趋零，矛盾。真正参数是 \(\lambda\epsilon\)。
\end{solution}

% 自拟：围绕本章已经建立的工具，改变条件或组合方法。
\begin{exercise}\label{b2:ex:21-new-7}
在小邻域内考虑 \(\Phi(x,y)=x\cdot y+|x|^2|y|^2\)。证明其满足\chapref{b2:ch:21}的振荡 \(L^2\) 估计。
\end{exercise}
\begin{solution}
混合 Hesse 矩阵 为 \(I+4xy^{\mathsf T}\)。当 \(|x||y|<1/8\) 时扰动算子范数小于 \(1/2\)，最小奇异值至少 \(1/2\)。振幅紧支撑于该邻域后，\cref{b2:thm:21-oscillatory-l2}给出 \(\|T_\lambda\|_{2\to2}\le C\lambda^{-d/2}\)。仅固定一个变量检查二阶 Hesse 矩阵 并不是这里所需的条件。
\end{solution}

% 自拟：围绕本章已经建立的工具，改变条件或组合方法。
\begin{optionalexercise}\label{b2:ex:21-new-8}
对半波相位 \(x\cdot\theta+t|\theta|-y\cdot\theta\)，计算典范关系，并说明 \(\theta=0\) 为什么应另行处理。
\end{optionalexercise}
\begin{solution}
驻点方程为 \(x+t\theta/|\theta|-y=0\)，两端余切向量均为 \(\theta\)，所以关系为 \((x,\theta;y,\theta)\)、\(y=x+t\theta/|\theta|\)、\(\theta\ne0\)。零点处 \(|\theta|\) 不光滑，齐次高频相位表示不适用；低频光滑截断产生的核可直接处理，不属于所讨论的高频典范关系。
\end{solution}
